\documentclass[12pt]{article}
\usepackage{graphicx} % for creating figures
\graphicspath{ {figures/} }

\usepackage{amsmath} % for math environments
\usepackage{amsfonts} % for math fonts like \mathbb
\usepackage{cancel} % for striking out math
\usepackage{bm} % for bold math



% stop typing \mathbb a thousand times 
\newcommand{\R}{\mathbb{R}}
\newcommand{\C}{\mathbb{C}}
\newcommand{\F}{\mathbb{F}}
\newcommand{\E}{\mathbb{E}}
\newcommand{\M}{\mathbb{M}}
\newcommand{\Z}{\mathbb{Z}}
\newcommand{\sphere}{\mathbb{S}}

\newcommand{\defeq}{\vcentcolon=} % for a prettier := sign
\newcommand{\eqdef}{=\vcentcolon} % for a prettier =: sign


% commands for bra-ket notation
\newcommand{\bra}[1]{\left\langle#1\right|}
\newcommand{\ket}[1]{\left|#1\right\rangle}
\newcommand{\braket}[2]{\left\langle #1 \middle| #2 \right\rangle}
\newcommand{\matrixel}[3]{\left\langle #1 \middle| #2 \middle| #3 \right\rangle}
\newcommand{\expectation}[1]{\left\langle #1 \right\rangle}


% vector stuff
\newcommand{\basis}{\bm{\hat{e}}}
\newcommand{\unit}[1]{\bm{\hat{#1}}}
\newcommand{\bvec}[1]{\bm{\vec{#1}}}
\newcommand{\del}{\bm{\vec{\nabla}}}

% curly velocity v 
\newcommand{\velocity}[1]{\vec{\mathcal{v}}_{{}_{#1}}}



\title{Dual Spaces}

\begin{document}
\maketitle

\section*{Linear Transformations as a Vector Space}
To begin, recall the definition of a linear transformation between two vector spaces $V$
and $W$, i.e. a mapping $\mathscr{L}:V\to W$ which satisfies
\begin{align}
  &\mathscr{L}\left(\vec{u}+\vec{v}\right) = \mathscr{L}\left( \vec{u} \right)+\mathscr{L}\left( \vec{v} \right) \\
  &\mathscr{L}\left( \alpha \vec{u} \right) = \alpha\mathscr{L}(\vec{u})
\end{align}
for any $\vec{u},\vec{v} \in V$ and $\alpha \in \R$. We note that operations on
the left hand side occur in $V$ while operations on the right hand side occur in $W$.

Having defined individual linear transformations, we now consider the set of
\textit{all} linear transformations between $V$ and $W$ which we denote
$\mathscr{L}(V,W)$. For any $\phi, \;\psi\in \mathscr{L}(V,W)$, we define
\begin{align}
  &(\phi+\psi)\vec{v} = \phi(\vec{v}) + \psi(\vec{v})\\
  &(\alpha \phi)(\vec{v}) = \alpha\phi(\vec{v})
\end{align}
Then ibt follows that $\mathscr{L}(V,W)$ is a vector space by linearity of each $\phi\in\mathscr{L}(V,W)$.
\end{document}






























